\begin{document}
\title{COMS3008A Assignment One -- Report}
\author{Your Full Name \\ Student Number: XXXXXXX}
\date{April 10, 2026}
\maketitle
\pagestyle{fancy}
\fancyhf{}
\fancyhead[R]{\thepage}
\fancyhead[L]{COMS3008A Assignment 1}
\pagenumbering{arabic}

% ────────────────────────────────────────────────
\section{Introduction}
% ────────────────────────────────────────────────

This report presents the parallelization of a Julia set fractal generator using OpenMP.
The Julia set is a fractal defined over the complex plane by iterating the function
$z_{n+1} = z_n^2 + c$, where $c = -0.7269 + 0.1889i$ is a fixed complex constant
\citep{julia1918memoire}.
A pixel at coordinates $(x,y)$ is considered part of the set if the iteration remains
bounded (does not exceed a magnitude threshold of 1000) within 300 iterations.
The image grid is $768 \times 768$ pixels.

Five parallel implementations of the compute-intensive pixel kernel were developed
using OpenMP shared-memory parallelism, and their performance was measured across
thread counts $\{1,2,4,6,8,10,12,14,16\}$.

% ────────────────────────────────────────────────
\section{Implementation Details}
% ────────────────────────────────────────────────

\subsection{Serial Baseline}

The serial implementation iterates over all rows $y \in [0, \text{DIM})$ and columns
$x \in [0, \text{DIM})$, evaluates \texttt{julia(x,y)} for each pixel, and writes
the result to a flat RGB buffer.
No synchronization is needed since each pixel is computed independently.
This independence is the key property that makes the problem \emph{embarrassingly
parallel} \citep{grama2003introduction}.

\subsection{Method (a): 1D Rowwise Parallel (Interleaved)}

In this method, assuming $T$ threads, thread $\textit{id}$ processes rows
$\textit{id},\ \textit{id}+T,\ \textit{id}+2T,\ \ldots$ in an interleaved (cyclic)
fashion.
For example with $T=4$: Thread~0 handles rows $0, 4, 8, \ldots$; Thread~1 handles
rows $1, 5, 9, \ldots$; and so on.
The decomposition is performed manually inside a \texttt{\#pragma omp parallel} region
using \texttt{omp\_get\_thread\_num()} and \texttt{omp\_get\_num\_threads()}.

\begin{lstlisting}[caption={1D Rowwise kernel (key excerpt)},label={lst:rowwise}]
#pragma omp parallel
{
    int id = omp_get_thread_num();
    int T  = omp_get_num_threads();
    for (int y = id; y < DIM; y += T) {
        for (int x = 0; x < DIM; x++) {
            int offset = x + y * DIM;
            int juliaValue = julia(x, y);
            ptr[offset*3+0] = 255 * juliaValue;
            ptr[offset*3+1] = 0;
            ptr[offset*3+2] = 0;
        }
    }
}
\end{lstlisting}

\subsection{Method (b): 1D Columnwise Parallel (Interleaved)}

This method mirrors Method~(a) but operates on columns instead of rows.
Thread $\textit{id}$ processes columns $\textit{id},\ \textit{id}+T,\
\textit{id}+2T,\ \ldots$
The outer loop iterates over $x$ (columns) and the inner loop over $y$ (rows).

\begin{lstlisting}[caption={1D Columnwise kernel (key excerpt)},label={lst:colwise}]
#pragma omp parallel
{
    int id = omp_get_thread_num();
    int T  = omp_get_num_threads();
    for (int x = id; x < DIM; x += T) {
        for (int y = 0; y < DIM; y++) {
            int offset = x + y * DIM;
            int juliaValue = julia(x, y);
            ptr[offset*3+0] = 255 * juliaValue;
            ptr[offset*3+1] = 0;
            ptr[offset*3+2] = 0;
        }
    }
}
\end{lstlisting}

\subsection{Method (c): 2D Row-Block Parallel (Contiguous Block)}

The $\text{DIM}$ rows are divided into $T$ contiguous blocks.
Thread $\textit{id}$ is assigned rows $[\textit{ystart},\ \textit{yend})$, where
$\textit{ystart} = \textit{id} \times \lfloor \text{DIM}/T \rfloor$ and
$\textit{yend} = (\textit{id}+1) \times \lfloor \text{DIM}/T \rfloor$, with the last
thread absorbing any remainder rows.

\begin{lstlisting}[caption={2D Row-Block kernel (key excerpt)},label={lst:rowblock}]
#pragma omp parallel
{
    int id             = omp_get_thread_num();
    int T              = omp_get_num_threads();
    int rows_per_thread = DIM / T;
    int ystart         = id * rows_per_thread;
    int yend           = (id == T-1) ? DIM : ystart + rows_per_thread;
    for (int y = ystart; y < yend; y++) {
        for (int x = 0; x < DIM; x++) { /* compute pixel */ }
    }
}
\end{lstlisting}

\subsection{Method (d): 2D Column-Block Parallel (Contiguous Block)}

Analogous to Method~(c) but the $\text{DIM}$ columns are divided into $T$ contiguous
blocks.
Thread $\textit{id}$ processes columns $[\textit{xstart},\ \textit{xend})$.
The outer loop iterates over rows $y$ and the inner loop is restricted to the
thread's column range.

\begin{lstlisting}[caption={2D Column-Block kernel (key excerpt)},label={lst:colblock}]
#pragma omp parallel
{
    int id              = omp_get_thread_num();
    int T               = omp_get_num_threads();
    int cols_per_thread = DIM / T;
    int xstart          = id * cols_per_thread;
    int xend            = (id == T-1) ? DIM : xstart + cols_per_thread;
    for (int y = 0; y < DIM; y++) {
        for (int x = xstart; x < xend; x++) { /* compute pixel */ }
    }
}
\end{lstlisting}

\subsection{Method (e): OpenMP \texttt{for} Construct}

The simplest approach: a single \texttt{\#pragma omp parallel for} directive with
static scheduling applied to the outer row loop.
OpenMP's runtime automatically divides the 768 iterations across the available threads
in contiguous blocks (static scheduling), closely resembling Method~(c) but with
far less manual code.

\begin{lstlisting}[caption={OpenMP for-construct kernel},label={lst:ompfor}]
#pragma omp parallel for schedule(static)
for (int y = 0; y < DIM; y++) {
    for (int x = 0; x < DIM; x++) {
        int offset = x + y * DIM;
        int juliaValue = julia(x, y);
        ptr[offset*3+0] = 255 * juliaValue;
        ptr[offset*3+1] = 0;
        ptr[offset*3+2] = 0;
    }
}
\end{lstlisting}

% ────────────────────────────────────────────────
\section{Performance Results}
% ────────────────────────────────────────────────

Table~\ref{tab:speedup} shows the measured speedup
($S = T_{\text{serial}} / T_{\text{parallel}}$) for each method across different
thread counts.
Figure~\ref{fig:speedup} plots the speedup curves.

\begin{table}[htb]
    \centering
    \caption{Speedup of each parallel method relative to the serial baseline.
             Results obtained on [FILL IN YOUR MACHINE SPECS].}
    \label{tab:speedup}
    \begin{tabular}{r|rrrrr}
        \toprule
        Threads & 1D Row & 1D Col & 2D Row-Block & 2D Col-Block & OMP For \\
        \midrule
        1  & FILL & FILL & FILL & FILL & FILL \\
        2  & FILL & FILL & FILL & FILL & FILL \\
        4  & FILL & FILL & FILL & FILL & FILL \\
        6  & FILL & FILL & FILL & FILL & FILL \\
        8  & FILL & FILL & FILL & FILL & FILL \\
        10 & FILL & FILL & FILL & FILL & FILL \\
        12 & FILL & FILL & FILL & FILL & FILL \\
        14 & FILL & FILL & FILL & FILL & FILL \\
        16 & FILL & FILL & FILL & FILL & FILL \\
        \bottomrule
    \end{tabular}
\end{table}

\begin{figure}[htb]
    \centering
    % Replace this with your actual pgfplots or includegraphics once you have results
    \fbox{\parbox{0.6\linewidth}{\centering
        \vspace{4cm}
        \textit{[Speedup line plot goes here.\\
        Paste your terminal.out numbers and ask Claude\\
        to regenerate this section with a pgfplots figure.]}
        \vspace{4cm}
    }}
    \caption{Speedup vs.\ number of threads for the five parallel implementations
             of the Julia set fractal kernel.}
    \label{fig:speedup}
\end{figure}

% ────────────────────────────────────────────────
\section{Discussion}
% ────────────────────────────────────────────────

\subsection{Question A: 1D Rowwise vs.\ 1D Columnwise}

Both methods assign the same number of pixels to each thread, so in terms of
\emph{computation load} they are balanced.
However, they differ significantly in \emph{memory access pattern}.

In Method~(a) (rowwise), each thread's inner loop walks along a row, accessing
consecutive memory locations: pixel offsets $x + y \cdot \text{DIM}$ are contiguous
for increasing $x$.
This access pattern exploits \emph{spatial locality} — consecutive pixels fall on the
same cache line — resulting in very few cache misses \citep{grama2003introduction}.

In Method~(b) (columnwise), the inner loop walks down a column.
Consecutive accesses are separated by $\text{DIM} \times 3 = 2304$ bytes (one full row
in the RGB buffer), which is likely to span many cache lines and cause frequent
\emph{cache misses}.
This makes Method~(b) a \emph{memory-bound} computation with noticeably lower speedup
than Method~(a), especially at higher thread counts where the memory bus becomes a
bottleneck.

\subsection{Question B: 2D Row-Block vs.\ 2D Column-Block}

The same cache locality argument applies.
Method~(c) (row-block) assigns each thread a contiguous range of rows; within those
rows the inner loop accesses memory sequentially — excellent spatial locality and good
cache performance.

Method~(d) (column-block) assigns each thread a contiguous range of columns.
Even though the column range for one thread is contiguous in the $x$-dimension, the
outer loop iterates over all 768 rows.
Within each row the thread writes only to its column slice
$[\textit{xstart}, \textit{xend})$, then the inner loop jumps to the next row.
The stride between successive writes by the same thread equals
$\text{DIM} \times 3$ bytes — again poor spatial locality.
Consequently, Method~(d) is expected to underperform Method~(c).

\subsection{Question C: Top Two Performers}

The two top-performing methods are \textbf{Method~(a) — 1D Rowwise} and
\textbf{Method~(e) — OpenMP \texttt{for} construct}.
Both achieve the highest speedup because they share the same fundamental
decomposition strategy: threads are assigned contiguous \emph{rows} (or interleaved
rows), meaning the inner loop always traverses consecutive memory addresses.
This yields high cache hit rates and low false-sharing risk.

Method~(e) with \texttt{schedule(static)} distributes rows in contiguous blocks
(identical to Method~(c)), but OpenMP's runtime also performs additional optimizations
such as minimizing thread-creation overhead through a pre-existing thread pool and
balancing the work division precisely.
Method~(a) achieves similarly good cache behavior through interleaving, which can
additionally improve load balance when the Julia set computation varies across rows
(some rows have many set-member pixels and therefore more iterations, while others
do not).

Both outperform the column-based methods because the underlying image buffer is
stored in \emph{row-major} order: iterating by rows accesses memory sequentially,
maximizing use of cache lines.

\subsection{Question D: Ideas for Further Improvement}

For \textbf{Method~(a) — 1D Rowwise (interleaved)}, potential further improvements
include:

\begin{itemize}
    \item \textbf{SIMD vectorization}: The inner $x$-loop operates independently on
    each pixel and involves only floating-point arithmetic on independent data.
    Adding compiler hints (e.g., \texttt{\#pragma omp simd}) or using AVX intrinsics
    could allow multiple pixels per row to be computed simultaneously within a single
    thread \citep{chapman2007using}.

    \item \textbf{Dynamic scheduling}: The Julia set is not uniform — some regions
    require the full 300 iterations while others exit early.
    Using \texttt{schedule(dynamic)} or \texttt{schedule(guided)} could better
    balance load across threads, particularly at lower thread counts.

    \item \textbf{Chunk-size tuning}: Increasing the chunk size in a dynamic schedule
    reduces scheduling overhead; tuning this parameter to the cache size of the target
    machine can improve performance.

    \item \textbf{Thread affinity}: Pinning threads to physical cores using
    \texttt{OMP\_PROC\_BIND=close} reduces cache migration overhead and ensures
    threads access memory local to their NUMA node \citep{trobec2018introduction}.
\end{itemize}

% ────────────────────────────────────────────────
\section{Conclusion}
% ────────────────────────────────────────────────

Five OpenMP parallel implementations of the Julia set fractal kernel were developed
and benchmarked.
The results confirm that memory access pattern is the dominant factor distinguishing
the methods: row-based decompositions benefit from sequential memory access and high
cache hit rates, while column-based methods suffer from strided access and increased
cache pressure.
The OpenMP \texttt{for} construct and the 1D rowwise interleaved method consistently
achieve the best speedup, closely approaching the theoretical linear speedup for the
embarrassingly parallel nature of the problem.

% ────────────────────────────────────────────────
\bibliographystyle{plainnat}
\bibliography{references}

\end{document}
